\section{The Original Concept}\label{sec:concept}

The 2013--2016 proposal describes a self-contained music-learning station
whose physical design and digital infrastructure work as one system.  The
following summary is drawn exclusively from the historical document
\cite{terzi_peeragogy_2016}; no contemporary reinterpretation is added here.

\subsection*{Physical structure}

The station is built around a drop-shaped enclosure housing the audio
diffusion system (Dolby Stereo Spectral Recording).  A digital keyboard
with MIDI input/output is the primary interaction device; during idle periods
the keyboard folds into a shelf, turning the station into an architectural
ornamental piece.  A foot pedal system provides sustain, soft, and resonance
control.  An OLED (or equivalent) screen serves as the visual interface,
switchable to a photo or video gallery when the station is not used for music.

\subsection*{Digital infrastructure}

A central computing unit handles audio processing, MIDI sequencing, video
playback, and web connectivity.  A webcam enables video transmission to
remote displays, supporting distance learning and remote performance.  The
system is designed to stream and reproduce audio--video concurrently,
allowing the performer to experience simultaneous sensory stimulation.

\subsection*{Educational model}

The proposal frames the station as an Open Educational Resource (OER)
delivery point: tutorial videos, music lessons, and a tutoring marketplace
are accessible through the unit.  The pedagogical backbone is peeragogical---a
collaborative, co-generated learning model documented in the companion text
\emph{The Peeragogy V3.3} and at peeragogy.org.  The intention is not to
replace a teacher but to create a low-friction access point for self-directed
musical exploration.

\subsection*{Context of the submission}

The document is filed as an invention disclosure---``a certain new and useful
Improvement in a Music Education field''---and is cast in the Free Technology
Guild seed-project format.  No explicit open-source or open-content license
is stated in the text.  The FTG context and the reference to
\emph{Peeragogy.org} suggest an intention to share the concept within a
collaborative community, but the original filing also names an inventor and
an application number, indicating a hybrid disclosure--sharing intent.
